@ID: OW24D1P1I
@BY: @refrain
@HINT: Gunakan sifat bahwa $(I_n)$ merupakan barisan interval menurun dan manfaatkan sifat Archimedes.

Misalkan

$$
I_n=\left(1,1+\frac1n\right).
$$

Terlihat bahwa

$$
I_1\supseteq I_2\supseteq I_3\supseteq\cdots.
$$

Maka

$$
\bigcup_{I_n\in\mathcal I}I_n
=I_1
=(1,2).
$$

Selanjutnya,

$$
\bigcap_{I_n\in\mathcal I}I_n
=\{x\in\mathbb R:x\in I_n,\ \forall n\in\mathbb N\}.
$$

Andaikan terdapat

$$
x\in\bigcap_{I_n\in\mathcal I}I_n.
$$

Karena $x>1$, berlaku $x-1>0$.

Dari sifat Archimedes, terdapat $n\in\mathbb N$ sehingga

$$
\frac1n<x-1.
$$

Akibatnya

$$
1+\frac1n<x,
$$

bertentangan dengan $x\in I_n$.

Jadi,

$$
\boxed{\bigcap_{I_n\in\mathcal I}I_n=\varnothing.}
$$


@ID: OW24D1P2I
@BY: @refrain
@HINT: Gunakan substitusi $x=\cos t$ dan identitas hasil kali ke jumlah.

Substitusi

$$
x=\cos t,\qquad t\in(0,\pi).
$$

Maka

$$
dx=-\sin t\,dt,\qquad
\sqrt{1-x^2}=\sin t,
$$

serta

$$
f_n(x)=\cos(nt),\qquad
f_m(x)=\cos(mt).
$$

Akibatnya,

$$
\int_{-1}^{1}\frac{f_n(x)f_m(x)}{\sqrt{1-x^2}}\,dx
=\int_0^\pi\cos(nt)\cos(mt)\,dt.
$$

Gunakan identitas

$$
\cos A\cos B
=\frac12\bigl(\cos(A+B)+\cos(A-B)\bigr),
$$

sehingga

$$
=\frac12\int_0^\pi
\left(\cos((m+n)t)+\cos((m-n)t)\right)\,dt
$$

$$
=\frac12\left[
\frac{\sin((m+n)t)}{m+n}
+\frac{\sin((m-n)t)}{m-n}
\right]_0^\pi
$$

$$
=0,
$$

karena $\sin(k\pi)=0$ untuk setiap $k\in\mathbb Z$.

Jadi,

$$
\boxed{0.}
$$

@ID: OW24D1P3I
@BY: @refrain
@HINT: Gunakan fakta bahwa pada grup siklis terdapat tepat $\varphi(d)$ elemen berorde $d$ untuk setiap pembagi $d$ dari orde grup.

Misalkan

$$
G=\langle a\rangle,
\qquad |G|=2024=2^3\cdot11\cdot23.
$$

Pada grup siklis, banyak elemen yang berorde $d$ adalah $\varphi(d)$ untuk setiap $d\mid2024$.

Pembagi ganjil dari $2024$ hanyalah

$$
1,\ 11,\ 23,\ 253.
$$

Maka banyak elemen yang berorde ganjil adalah

$$
\varphi(1)+\varphi(11)+\varphi(23)+\varphi(253)
=1+10+22+220
=253.
$$

Jadi jawabannya adalah

$$
\boxed{253.}
$$
@ID: OW24D1P4I
@BY: @aduhmath
@HINT: Gunakan fakta bahwa konjugat-konjugat aljabar dari $\sqrt3-\sqrt5$ juga merupakan akar.

Karena

$$
\sqrt3-\sqrt5
$$

merupakan akar polinom berkoefisien bilangan bulat, maka

$$
\sqrt3+\sqrt5,\quad
-\sqrt3+\sqrt5,\quad
-\sqrt3-\sqrt5
$$

juga merupakan akar.

Akibatnya,

$$
\begin{aligned}
f(x)
&=(x+\sqrt3+\sqrt5)(x+\sqrt3-\sqrt5)
(x-\sqrt3+\sqrt5)(x-\sqrt3-\sqrt5)\\
&=\bigl((x+\sqrt3)^2-5\bigr)
\bigl((x-\sqrt3)^2-5\bigr)\\
&=(x^2-2+2x\sqrt3)(x^2-2-2x\sqrt3)\\
&=(x^2-2)^2-12x^2\\
&=x^4-16x^2+4.
\end{aligned}
$$

Maka

$$
a=0,\qquad
b=-16,\qquad
c=0,\qquad
d=4,
$$

sehingga

$$
a+b+c+d
=0-16+0+4
=-12.
$$

Jadi jawabannya adalah

$$
\boxed{-12.}
$$

@ID: OW24D1P5I
@BY: @hniv
@HINT: Nyatakan banyak kelereng yang diambil dari setiap warna sebagai solusi bilangan bulat taknegatif, lalu gunakan metode stars and bars.

Misalkan $x_i$ menyatakan banyaknya kelereng warna ke-$i$ yang diambil, untuk $i=1,2,\ldots,5$.

Karena setiap warna terambil paling sedikit satu kali, berlaku
$x_1+x_2+x_3+x_4+x_5=9$ dan $x_i\ge1$.

Substitusi $y_i=x_i-1$, sehingga $y_i\ge0$ dan

$$
y_1+y_2+y_3+y_4+y_5=4.
$$

Karena $y_i\le4<9$, tidak ada batasan atas yang aktif. Dengan metode stars and bars, banyak solusi adalah

$$
\binom{5+4-1}{5-1}
=\binom84
=70.
$$

Jadi jawabannya adalah $\boxed{70}$.


@ID: OW24D1P1U
@BY: @fofo
@HINT: Tunjukkan bahwa fungsi iterasi merupakan kontraksi, lalu terapkan Teorema Titik Tetap Banach.

Karena $\sqrt b<a$, diperoleh $a^2-b>0$.

Definisikan fungsi $f:\mathbb R_{\ge0}\to\mathbb R_{\ge0}$ dengan

$$
f(x)=\frac{ax+b}{x+a}
=a-\frac{a^2-b}{x+a}.
$$

Untuk setiap $x,y\ge0$,

$$
\begin{aligned}
|f(x)-f(y)|
&=\left|\left(a-\frac{a^2-b}{x+a}\right)
-\left(a-\frac{a^2-b}{y+a}\right)\right|\\
&=(a^2-b)\left|\frac1{x+a}-\frac1{y+a}\right|\\
&=(a^2-b)\frac{|x-y|}{(x+a)(y+a)}\\
&\le\frac{a^2-b}{a^2}|x-y|\\
&=\left(1-\frac{b}{a^2}\right)|x-y|.
\end{aligned}
$$

Karena

$$
0<1-\frac{b}{a^2}<1,
$$

maka $f$ merupakan fungsi kontraksi.

Misalkan $f^{(n)}$ menyatakan komposisi $f$ sebanyak $n$ kali. Berdasarkan Teorema Titik Tetap Banach, terdapat tepat satu $c\ge0$ sehingga $f(c)=c$, dan untuk setiap $x\ge0$, barisan $\bigl(f^{(n)}(x)\bigr)$ konvergen ke $c$.

Karena $x_n=f(x_{n-1})$, diperoleh $x_n=f^{(n)}(x_0)$ untuk setiap $n\ge1$. Jadi, $(x_n)$ konvergen ke $c$.

Selanjutnya,

$$
\frac{ac+b}{c+a}=c
\iff
ac+b=c^2+ac
\iff
c^2=b.
$$

Karena $c\ge0$, diperoleh

$$
c=\sqrt b.
$$

Dengan demikian,

$$
\boxed{\lim_{n\to\infty}x_n=\sqrt b.}
$$

@ID: OW24D1P2U
@BY: @aduhmath
@HINT: Gunakan Teorema Nilai Ekstrem pada fungsi $|f|$.

Karena $f$ kontinu, maka $|f|$ juga kontinu pada $[a,b]$.

Karena $[a,b]$ kompak, berdasarkan Teorema Nilai Ekstrem terdapat $x_0\in[a,b]$ sehingga

$$
|f(x_0)|=\min_{x\in[a,b]}|f(x)|.
$$

Dari hipotesis, untuk $x=x_0$ terdapat $p\in[a,b]$ sehingga

$$
|f(p)|\le\frac{2023}{2024}|f(x_0)|.
$$

Karena $|f(x_0)|$ merupakan nilai minimum dari $|f|$, diperoleh

$$
|f(x_0)|
\le
|f(p)|
\le
\frac{2023}{2024}|f(x_0)|.
$$

Maka

$$
\frac1{2024}|f(x_0)|\le0.
$$

Akibatnya,

$$
|f(x_0)|=0,
$$

sehingga

$$
f(x_0)=0.
$$

Dengan mengambil $c=x_0$, diperoleh

$$
\boxed{f(c)=0.}
$$

@ID: OW24D1P3U
@BY: @refrain
@HINT: Misalkan $(a,b)$ pembagi nol dan analisis berdasarkan komponen pertama dari hasil perkalian.

Misalkan $(a,b)\ne(0,0)$ merupakan pembagi nol. Maka terdapat $(c,d)\ne(0,0)$ sehingga

$$
(a,b)\cdot(c,d)=(0,0).
$$

Artinya,

$$
ac\equiv0\pmod{1013},
\qquad
ad+bc+bd\equiv0\pmod{1013}.
$$

Karena $1013$ prima, dari $ac\equiv0\pmod{1013}$ diperoleh $a=0$ atau $c=0$.

**Kasus 1.** $a=0$.

Karena $(a,b)\ne(0,0)$, maka $b\ne0$. Persamaan kedua menjadi

$$
bc+bd=b(c+d)\equiv0\pmod{1013}.
$$

Karena $b\ne0$, diperoleh

$$
c+d\equiv0\pmod{1013}.
$$

Untuk setiap $b\in\{1,2,\ldots,1012\}$, dapat dipilih misalnya $(c,d)=(1,1012)$ sehingga syarat di atas terpenuhi. Jadi terdapat $1012$ pembagi nol pada kasus ini.

**Kasus 2.** $c=0$.

Karena $(c,d)\ne(0,0)$, maka $d\ne0$. Persamaan kedua menjadi

$$
ad+bd=d(a+b)\equiv0\pmod{1013}.
$$

Karena $d\ne0$, diperoleh

$$
a+b\equiv0\pmod{1013}.
$$

Maka pembagi nol yang mungkin adalah

$$
(1,1012),(2,1011),\ldots,(1012,1),
$$

yang berjumlah $1012$ buah.

Dengan demikian, banyak elemen taknol di $R$ yang merupakan pembagi nol adalah

$$
1012+1012=2024.
$$

Jadi terdapat tepat

$$
\boxed{2024}
$$

elemen taknol di $R$ yang merupakan pembagi nol.


@ID: OW24D1P4U
@BY: @refrain
@HINT: Untuk (a) cukup berikan konstruksi. Untuk (b), jumlahkan seluruh elemen pada kedua himpunan.

**a)**

Ambil

$$
g_i=\overline{i}\in\mathbb Z_7,
\qquad i=1,2,\ldots,7.
$$

Maka

$$
\{g_1+g_2,\;g_2+g_3,\;\ldots,\;g_6+g_7,\;g_7+g_1\}
=\mathbb Z_7.
$$

Jadi $(\mathbb Z_7,+)$ merupakan grup yang rapi.

**b)**

Andaikan $(\mathbb Z_n,+)$ rapi untuk suatu $n$ genap. Maka terdapat $g_1,g_2,\ldots,g_n$ berbeda sehingga

$$
\{g_1+g_2,\;g_2+g_3,\;\ldots,\;g_n+g_1\}
=\mathbb Z_n.
$$

Menjumlahkan semua elemen pada kedua himpunan diperoleh

$$
\sum_{i=1}^n(g_i+g_{i+1})
=
2\sum_{i=1}^n g_i
=
\sum_{i=1}^n\overline{i}
\pmod n,
$$

dengan $g_{n+1}=g_1$.

Karena $g_1,\ldots,g_n$ merupakan permutasi dari elemen-elemen $\mathbb Z_n$,

$$
\sum_{i=1}^n g_i
=
\sum_{i=1}^n\overline{i}.
$$

Akibatnya,

$$
2\sum_{i=1}^n\overline{i}
=
\sum_{i=1}^n\overline{i}
\pmod n,
$$

sehingga

$$
\sum_{i=1}^n\overline{i}
=
0
\pmod n.
$$

Namun,

$$
\sum_{i=1}^n\overline{i}
=
\frac{n(n+1)}2,
$$

sedangkan untuk $n$ genap,

$$
\frac{n(n+1)}2
\not\equiv0
\pmod n,
$$

karena $n+1$ ganjil.

Kontradiksi.

Jadi, untuk setiap $n$ genap, $(\mathbb Z_n,+)$ tidak rapi.


@ID: OW24D1P5U
@BY: @hniv
@HINT: Gunakan prinsip Pigeonhole dengan mempartisi himpunan menjadi $n$ blok berukuran $k$.

Partisi himpunan $\{1,2,\ldots,kn\}$ menjadi $n$ himpunan

$$
H_i=\{k(i-1)+1,\;k(i-1)+2,\;\ldots,\;ki\},
\qquad i=1,2,\ldots,n.
$$

Setiap $H_i$ berisi tepat $k$ bilangan.

Karena dipilih $n+1$ bilangan, sedangkan hanya terdapat $n$ himpunan partisi, berdasarkan Prinsip Pigeonhole terdapat dua bilangan yang berada pada himpunan $H_i$ yang sama.

Selisih maksimum dua bilangan di dalam $H_i$ adalah

$$
ki-\bigl(k(i-1)+1\bigr)=k-1.
$$

Dengan demikian, selalu terdapat dua bilangan yang selisihnya paling banyak $k-1$.

@ID: OW24D2P1I
@BY: @refrain
@HINT: Gunakan fakta bahwa $z^{20}=1$ sehingga $|z|=1$.

Karena $z^{20}=1$ dan $z^{24}\in\mathbb R$, berlaku

$$
z^4=\frac{z^{24}}{z^{20}}=z^{24}\in\mathbb R.
$$

Karena $z^{20}=1$, maka $|z|=1$, sehingga

$$
z^4\in\{-1,1\}.
$$

Apabila $z^4=-1$, maka

$$
z^{20}=(z^4)^5=-1,
$$

bertentangan dengan $z^{20}=1$.

Jadi haruslah

$$
z^4=1.
$$

Banyak akar dari $z^4=1$ adalah $4$.

Jadi jawabannya adalah

$$
\boxed{4.}
$$

@ID: OW24D2P2I
@BY: @wildabandon
@HINT: Tuliskan $z=x+iy$, lalu gunakan persamaan Cauchy--Riemann.

Tuliskan $z=x+iy$ dengan $x,y\in\mathbb R$.

Maka

$$
f(x+iy)
=(x-iy)e^{-x^2-y^2}.
$$

Dengan demikian,

$$
u(x,y)=xe^{-x^2-y^2},
\qquad
v(x,y)=-ye^{-x^2-y^2}.
$$

Turunan parsialnya memenuhi

$$
u_x=e^{-x^2-y^2}-2x^2e^{-x^2-y^2},
\qquad
u_y=-2xye^{-x^2-y^2},
$$

$$
v_x=2xye^{-x^2-y^2},
\qquad
v_y=-e^{-x^2-y^2}+2y^2e^{-x^2-y^2}.
$$

Di titik $(0,1)$ diperoleh

$$
u_x(0,1)=\frac1e,\qquad
u_y(0,1)=0,
$$

$$
v_x(0,1)=0,\qquad
v_y(0,1)=\frac1e.
$$

Karena persamaan Cauchy--Riemann berlaku di sekitar $(0,1)$, maka

$$
f'(i)
=u_x(0,1)+iv_x(0,1)
=\frac1e.
$$

Jadi jawabannya adalah

$$
\boxed{\frac1e.}
$$

@ID: OW24D2P3I
@BY: @hniv
@HINT: Hitung bentuk $T_k^m$ untuk setiap $m$.

Untuk setiap polinomial

$$
p(x)=a_0+a_1x+a_2x^2+a_3x^3+a_4x^4,
$$

diperoleh

$$
T_k(p)=a_0+a_1^kx^2+a_2x^4,
$$

$$
T_k^2(p)=a_0+a_1^kx^4,
$$

$$
T_k^3(p)=a_0,
$$

dan

$$
T_k^4(p)=a_0.
$$

Jelas $T_k^3$ dan $T_k^4$ merupakan pemetaan linear untuk setiap $k$.

Selanjutnya, $T_k$ dan $T_k^2$ linear jika dan hanya jika $k=1$.

Maka pasangan $(k,m)$ yang memenuhi adalah

$$
(1,1),(1,2),(1,3),(1,4),
$$

$$
(2,3),(2,4),
$$

$$
(3,3),(3,4),
$$

$$
(4,3),(4,4),
$$

yang berjumlah $10$.

Jadi jawabannya adalah

$$
\boxed{10.}
$$

@ID: OW24D2P4I
@BY: @hniv
@HINT: Dimensi subruang terkecil yang memuat $U$ dan $V$ sama dengan dimensi $\operatorname{span}(U\cup V)$.

Karena subruang terkecil yang memuat $U$ dan $V$ adalah $\operatorname{span}(U\cup V)$, cukup mencari

$$
\dim(\operatorname{span}(U\cup V)).
$$

Bentuk matriks $A$ yang kolom-kolomnya adalah vektor-vektor pada $U\cup V$. Maka

$$
\dim(\operatorname{span}(U\cup V))
=\operatorname{rank}(A).
$$

Dengan mengambil transpose, diperoleh

$$
A^T=
\begin{bmatrix}
1&0&0&0&0&0&0&0&0&0\\
0&1&1&0&0&0&0&0&0&0\\
0&0&0&1&1&1&0&0&0&0\\
0&0&0&0&0&0&1&1&1&1\\
1&1&0&0&0&0&0&0&0&0\\
0&0&1&1&0&0&0&0&0&0\\
0&0&0&0&1&1&0&0&0&0\\
0&0&0&0&0&0&1&1&0&0\\
0&0&0&0&0&0&0&0&1&1
\end{bmatrix}.
$$

Dengan melakukan operasi baris elementer, matriks tersebut ekuivalen dengan matriks yang memiliki $7$ baris taknol.

Akibatnya,

$$
\operatorname{rank}(A)
=\operatorname{rank}(A^T)
=7.
$$

Jadi,

$$
\dim(\operatorname{span}(U\cup V))=7.
$$

Dengan demikian, dimensi subruang terkecil yang memuat $U$ dan $V$ adalah

$$
\boxed{7.}
$$

@ID: OW24D2P5I
@BY: @wildabandon
@HINT: Gunakan jumlah total skor seluruh peserta.

Banyak pertandingan yang berlangsung adalah

$$
\binom{10}{2}=45.
$$

Karena pada setiap pertandingan jumlah skor kedua peserta selalu $2$, maka jumlah seluruh skor peserta adalah

$$
2\cdot45=90.
$$

Misalkan skor peserta dengan skor terendah paling sedikit $5$. Karena seluruh skor berbeda, jumlah skor seluruh peserta paling sedikit

$$
5+6+\cdots+14
=\frac{(5+14)\cdot10}{2}
=95,
$$

bertentangan dengan fakta bahwa jumlah seluruh skor adalah $90$.

Jadi skor peserta dengan skor terendah paling besar adalah $4$.

Karena satu kemenangan bernilai $2$ poin, peserta dengan skor $4$ paling banyak dapat memperoleh

$$
\frac42=2
$$

kemenangan.

Batas ini dapat dicapai (contoh konfigurasi diserahkan pada pembaca).

$$
\boxed{2}.
$$


@ID: OW24D2P1U
@BY: @refrain
@HINT: Gunakan ketaksamaan segitiga dan ketaksamaan segitiga balik.

Dengan ketaksamaan segitiga balik,

$$
||z|-1|
\le
|1+z|
\le
|z|+1.
$$

Karena $|z|=1$, diperoleh

$$
0\le |1+z|\le2.
$$

Demikian pula,

$$
||2z|-1|
\le
|1+2z|
\le
|2z|+1.
$$

Karena $|2z|=2$, diperoleh

$$
1\le |1+2z|\le3.
$$

Menjumlahkan kedua ketaksamaan tersebut menghasilkan

$$
1=0+1
\le
|1+z|+|1+2z|
\le
2+3=5.
$$

Jadi,

$$
\boxed{1\le |1+z|+|1+2z|\le5.}
$$

@ID: OW24D2P2U
@BY: @fofo (direvisi @refrain)
@HINT: Carilah daerah tempat $\operatorname{Re}(z)\to-\infty$ ketika $|z|\to\infty$.

Untuk setiap bilangan real $a,b_1,b_2$ dengan $b_1<b_2$, definisikan

$$
D(a,b_1,b_2)
=
\{z\in\mathbb C:\operatorname{Re}(z)<a,\ b_1<\operatorname{Im}(z)<b_2\}.
$$

**Lemma.** Untuk setiap $z\in D(a,b_1,b_2)$,

$$
\lim_{\substack{|z|\to\infty\\z\in D(a,b_1,b_2)}}
\operatorname{Re}(z)
=
-\infty.
$$

**Bukti.**

Karena $z\in D(a,b_1,b_2)$,

$$
|\operatorname{Im}(z)|
\le
\max\{|b_1|,|b_2|\}
=:B.
$$

Selain itu,

$$
\operatorname{Re}(z)<a.
$$

Dari

$$
|z|
\le
|\operatorname{Re}(z)|
+
|\operatorname{Im}(z)|
\le
|\operatorname{Re}(z)|+B,
$$

diperoleh $|\operatorname{Re}(z)|\to\infty$ ketika $|z|\to\infty$.

Karena $\operatorname{Re}(z)<a$, haruslah

$$
\operatorname{Re}(z)\to-\infty.
$$

Selanjutnya,

$$
|e^z|
=
e^{\operatorname{Re}(z)}.
$$

Maka, berdasarkan lemma,

$$
\lim_{\substack{|z|\to\infty\\z\in D(a,b_1,b_2)}}
|e^z|
=
\lim_{\operatorname{Re}(z)\to-\infty}
e^{\operatorname{Re}(z)}
=
0.
$$

Akibatnya,

$$
\lim_{\substack{|z|\to\infty\\z\in D(a,b_1,b_2)}}e^z
=
0.
$$

Dengan demikian, salah satu daerah yang memenuhi adalah

$$
\boxed{
D=
\{z\in\mathbb C:\operatorname{Re}(z)<a,\ b_1<\operatorname{Im}(z)<b_2\},
}
$$

dengan $a,b_1,b_2\in\mathbb R$ dan $b_1<b_2$ sebarang.

@ID: OW24D2P3U
@BY: @refrain ft. @wili
@HINT: Tunjukkan bahwa $A$ memiliki tiga nilai eigen yang berbeda.

Polinom karakteristik $A$ adalah

$$
\det(\lambda I-A)
=
\begin{vmatrix}
\lambda&-a&0\\
-a&\lambda-b&-a\\
0&-a&\lambda
\end{vmatrix}.
$$

Ekspansi terhadap baris pertama menghasilkan

$$
\det(\lambda I-A)
=
\lambda(\lambda(\lambda-b)-a^2)-a^2\lambda
=
\lambda(\lambda^2-b\lambda-2a^2).
$$

Jadi nilai-nilai eigennya adalah

$$
\lambda_1=0,
$$

dan

$$
\lambda_{2,3}
=
\frac{b\pm\sqrt{b^2+8a^2}}2.
$$

Karena $a\ne0$,

$$
b^2+8a^2>b^2,
$$

sehingga

$$
\sqrt{b^2+8a^2}>|b|.
$$

Akibatnya,

$$
\lambda_2\ne\lambda_3,
\qquad
\lambda_2\ne0,
\qquad
\lambda_3\ne0.
$$

Maka $A$ memiliki tiga nilai eigen yang saling berbeda. Oleh karena itu, $A$ memiliki tiga vektor eigen yang bebas linear.

Dengan demikian, $A$ dapat didiagonalkan.

@ID: OW24D2P4U
@BY: @RM (diketik @refrain)
@HINT: Agar sistem tidak memiliki solusi, $A$ harus singular.

Agar terdapat $b$ sehingga persamaan

$$
A^2x=Ab
$$

tidak memiliki solusi, matriks $A$ harus singular.

Karena

$$
\det(A)=20q-24p,
$$

maka harus berlaku

$$
20q-24p=0,
$$

atau

$$
q=\frac65p.
$$

Dengan syarat tersebut,

$$
A=
\begin{pmatrix}
20&24\\
p&\frac65p
\end{pmatrix},
$$

dan kedua kolom $A$ saling bergantung.

Untuk setiap $x=(x_1,x_2)^T$, berlaku

$$
Ax
=
k_1
\begin{pmatrix}
20\\
p
\end{pmatrix},
$$

dengan

$$
k_1=x_1+\frac65x_2.
$$

Demikian pula,

$$
Ab
=
k_2
\begin{pmatrix}
20\\
p
\end{pmatrix},
$$

untuk suatu $k_2\in\mathbb R$.

Selanjutnya,

$$
A^2x=A(Ax)
=
k_1A
\begin{pmatrix}
20\\
p
\end{pmatrix}
=
k_1\left(20+\frac65p\right)
\begin{pmatrix}
20\\
p
\end{pmatrix}.
$$

Maka persamaan $A^2x=Ab$ ekuivalen dengan

$$
k_1\left(20+\frac65p\right)=k_2.
$$

Jika

$$
20+\frac65p\ne0,
$$

maka untuk setiap $k_2$ selalu dapat dipilih

$$
k_1=\frac{k_2}{20+\frac65p},
$$

sehingga persamaan selalu mempunyai solusi.

Jadi satu-satunya kemungkinan adalah

$$
20+\frac65p=0,
$$

yakni

$$
p=-\frac{50}{3},
\qquad
q=-20.
$$

Untuk pasangan ini,

$$
A^2=0,
$$

sedangkan $\operatorname{rank}(A)=1$, sehingga terdapat $b\notin\ker(A)$.

Akibatnya,

$$
Ab\ne0=A^2x
$$

untuk setiap $x$, sehingga persamaan tidak memiliki solusi.

Dengan demikian, satu-satunya pasangan yang memenuhi adalah

$$
\boxed{\left(-\frac{50}{3},-20\right).}
$$

@ID: OW24D2P5U
@BY: @wili (diketik @refrain)
@HINT: Hitung semua pasangan $(a,b)$, lalu identifikasi pasangan yang merepresentasikan garis yang sama.

Penyelesaian dibagi menjadi tiga kasus.

**Kasus 1.** $a=0$.

Persamaan menjadi $by=0$ dengan $b\in\{1,2,3,6,7\}$.

Semua pilihan $b$ merepresentasikan garis yang sama, yaitu

$$
y=0.
$$

Jadi terdapat $1$ garis pada kasus ini.

**Kasus 2.** $b=0$.

Secara analog, semua pilihan $a$ merepresentasikan garis

$$
x=0.
$$

Jadi terdapat $1$ garis pada kasus ini.

**Kasus 3.** $a,b\neq0$.

Banyak pasangan terurut $(a,b)$ dengan $a,b\in\{1,2,3,6,7\}$ dan $a\neq b$ adalah

$$
P(5,2)=\frac{5!}{(5-2)!}=20.
$$

Beberapa pasangan berbeda merepresentasikan garis yang sama, yaitu

- $(1,2)$ dan $(3,6)$,
- $(2,1)$ dan $(6,3)$,
- $(1,3)$ dan $(2,6)$,
- $(3,1)$ dan $(6,2)$.

Keempat pasangan tersebut masing-masing mengurangi satu garis yang dihitung ganda.

Maka banyak garis berbeda pada kasus ini adalah

$$
20-4=16.
$$

Dengan demikian, banyak garis lurus berbeda yang dapat dibentuk adalah

$$
1+1+16=18.
$$

Jadi jawabannya adalah

$$
\boxed{18}.
$$
